\section{Introduction}\label{sec:intro}

Neurosymbolic systems combine neural perception with symbolic reasoning, but the two halves run on different machines. Deep learning frameworks---PyTorch, JAX---execute dense tensor computation on the GPU through optimized CUDA kernels. Logic engines---Datalog evaluators, probabilistic-logic systems such as ProbLog~\cite{problog2}, answer-set solvers---are CPU-bound interpreters and compilers. Every system that couples the two, from DeepProbLog~\cite{deepproblog} to NeurASP~\cite{neurasp}, therefore straddles the PCIe bus: each training iteration ships a network's outputs to a CPU logic engine, materializes query results or proofs, and pulls gradients back to update parameters. At scale---millions of ground atoms, thousands of steps---these host--device transfers, not the inference or learning itself, dominate wall-clock time and bandwidth.

The transfer wall is one half of the problem; partial coverage is the other. Existing neurosymbolic frameworks bridge a neural network to a \emph{single} symbolic backend---DeepProbLog to probabilistic logic, NeurASP to answer-set programming---and run that backend on the CPU. GPU logic efforts, conversely, accelerate one paradigm in isolation: deterministic Datalog~\cite{gpulog,vflog} or SAT~\cite{parafrost}, with no neural interface. No system makes the \emph{whole} symbolic spectrum---deterministic, probabilistic, and epistemic reasoning---both GPU-resident and differentiably coupled to neural networks.

xlog closes both gaps. It is a GPU-native logic programming language whose compiler and runtime treat a neural network as an ordinary predicate: the network's softmax output becomes a distribution over weighted facts that flow into whichever reasoning mode a program uses, and gradients flow back, all without leaving the device. Three mechanisms make this integration both universal and sound. A \emph{device-resident symbolic substrate} keeps fact stores, deltas, circuits, solver state, and gradient buffers in GPU memory throughout evaluation, so the symbolic half never crosses the bus. \emph{Verified knowledge compilation} turns a network's outputs into a differentiable arithmetic circuit and proves, on-device, that the circuit is logically equivalent to its source formula---so the neural--symbolic bridge is exact and machine-checked, not merely assumed. \emph{Zero-copy differentiable coupling} exports the resulting gradients through DLPack~\cite{dlpack} directly into PyTorch's autograd graph. The system is implemented in Rust with custom CUDA kernels organized into a layered crate architecture; host involvement is limited to orchestration, I/O, and compilation.

The principal contributions of this paper are:

\begin{itemize}
  \item \textbf{A universal neurosymbolic integration model.} A neural network is a predicate, and the same uniform, transfer-free mechanism connects it to deterministic, probabilistic, and epistemic reasoning---rather than a bespoke bridge to one backend (\Cref{sec:arch,sec:nesy}).
  \item \textbf{A GPU-resident symbolic substrate.} Semi-naive Datalog evaluation with stratified negation and aggregation runs as custom CUDA relational kernels, extended with a worst-case-optimal join (WCOJ) subsystem that avoids the intermediate-relation blowup of binary-join plans on skewed cyclic queries (a measured $27.96\times$ geometric-mean ablation, \Cref{sec:datalog}).
  \item \textbf{GPU-native verified knowledge compilation.} A provenance$\to$CNF$\to$D4$\to$circuit pipeline runs entirely on device, and an on-GPU CDCL solver proves each compiled circuit equivalent to its source formula. To our knowledge this is the first end-to-end GPU-resident knowledge-compilation pipeline with on-device equivalence verification, and it is the paper's central technical result (\Cref{sec:prob}).
  \item \textbf{End-to-end differentiable training with zero-copy interop.} Compiled circuits are cached across iterations and evaluated as level-parallel forward/backward kernels, with gradients exported zero-copy via DLPack and Arrow; circuit caching yields a measured $2.74\times$ training speedup. Term embeddings and differentiable ILP ride the same device-resident path (\Cref{sec:nesy}).
  \item \textbf{Full-spectrum reach.} Epistemic reasoning over world views (\texttt{know}/\texttt{possible} under founded and Gelfond-style semantics) and well-founded semantics execute on the same substrate, reusing the CDCL, join, and circuit machinery rather than a separate engine (\Cref{sec:epistemic}).
\end{itemize}

The remainder of this paper is organized as follows. \Cref{sec:arch} presents the integration model and system architecture. \Cref{sec:lang} describes the xlog language. \Cref{sec:datalog} details the GPU-resident symbolic substrate. \Cref{sec:prob} presents verified knowledge compilation, the central contribution. \Cref{sec:nesy} ties the pieces together as end-to-end neurosymbolic learning. \Cref{sec:epistemic} extends the substrate to epistemic and other reasoning modes. \Cref{sec:eval} reports evaluation results, \Cref{sec:related} surveys related work, and \Cref{sec:limits} discusses limitations and future work.
